\makeabstract
{
Deriving Temperatures for Substellar Objects Within 20 Parsecs Using New Dynamical Masses 
}
{
Simone Lilavois (1), Ryan Purviance, (1), Mark R. Giovinazzi, PhD (1), Daniella Bardalez Gagliuffi, PhD (1)
}
{
\textsuperscript{1}Amherst College, Amherst, MA, USA
}
{
To deepen our understanding of the solar system and Earth’s place in the galaxy, we must first obtain a sample of reliable spectra to investigate trends in the chemical compositions of substellar atmospheres. This can be achieved with direct imaging, which requires high-precision measurements of dynamical masses and orbital parameters. Using orvara, a Markov Chain Monte Carlo approach with state-of-the-art statistical techniques for efficient orbit solving, we derive, for the first time, complete system parameters for 29 objects in 25 systems: 21 stars, 2 brown dwarfs, and 6 planets.
}
{
By combining 5,000+ radial velocity measurements, 1,000+ relative astrometric observations from the Washington Double Star Catalog, and absolute astrometry from the Gaia DR3 version of the Hipparcos–Gaia Catalog of Accelerations, we calculate companion masses dynamically, improving upon previous model-dependent masses, as well as minimum masses that have inclination degeneracies. With these true masses and orbital parameters, notably semi-major axes, eccentricities, and inclinations, we constrain the three-dimensional architectures of stellar and substellar systems. Our new dynamical masses, paired with known host star ages, are used as input to substellar evolutionary models to derive the first-ever temperatures for five objects.
}
{
Using substellar evolutionary models, we calculate new temperatures for 2 brown dwarfs: HIP 1292B (1330.3 K) and HIP 10812B (1119.4 K), as well as 3 exoplanets: HIP 1292Ab (381.8 K), HIP 71395b (249.6 K), and HIP 71395c (298.8 K). This establishes each as an exciting candidate for atmospheric characterization with JWST, which has a threshold of \textasciitilde{}300 K for optimal direct imaging. 
}
{
Our robust sample represents the newest and highest precision constraints on systems in the solar neighborhood and will set the stage for the upcoming Nancy Grace Roman Space Telescope and Habitable Worlds Observatory, which will be capable of directly imaging Earth-like exoplanets.
}

\newpage
